\documentclass[12pt,a4paper]{article}

\usepackage[utf8]{inputenc}
\usepackage[T1]{fontenc}
\usepackage[margin=2.5cm]{geometry}
\usepackage{hyperref}
\usepackage{microtype}
\usepackage{parskip}
\usepackage{lmodern}
\usepackage{amsmath}

\hypersetup{colorlinks=true, urlcolor=blue!60!black}

\pagestyle{empty}

\begin{document}

\begin{flushright}
Kundai Farai Sachikonye \\
Auf der Lichtung 19, 82178 Puchheim, Germany \\
\href{https://github.com/fullscreen-triangle}{github.com/fullscreen-triangle} \\
\texttt{kundai.sachikonye@bitspark.com} \\
(+49) 1626386344 \\[6pt]
22 July 2026
\end{flushright}

\vspace{1em}

Dr.\ med.\ Felix Busch \\
Chair of Medical Informatics \\
Institute for AI and Informatics in Medicine (AIIM) \\
TUM School of Medicine and Health \\
Ismaninger Str.\ 22, 81675 Munich

\vspace{1.5em}

\textbf{Re: Doctoral position, xMDT-HPB --- explainable multimodal AI for
hepato-pancreato-biliary oncology}

\vspace{1em}

Dear Dr.\ Busch,

I am applying for the doctoral position in the xMDT-HPB project.
I want to do this PhD, and to do it inside a structured clinical-informatics
group, for a specific reason: I have spent several years building explainable
AI systems for cancer on my own, and I have reached the limit of what
independent work can teach me about the part that matters most here ---
making those systems trustworthy to clinicians inside a real multidisciplinary
tumour board, against real clinical data and real standards.
That is the education a structured doctorate at AIIM offers and independent
work cannot, and it is the reason I am applying rather than continuing alone.

\textbf{The project's central word is the one I have worked on hardest.}
xMDT-HPB is built around AI that clinicians can \emph{understand and trust}.
Most work on this responds to the black-box problem with post-hoc explanation
--- SHAP values, attention maps, saliency --- which documents a model's behaviour
without making the model itself auditable.
I have spent my independent research building diagnostic systems that are
auditable in a stronger sense.
The \href{https://syndrome-seven.vercel.app/tools}{Syndrome} framework performs
MHC binding prediction (AUC\,$0.81$, validated on 2{,}847 IEDB peptides),
neoantigen identification, and cancer immune-escape scoring using three-channel
discrete Fourier transforms on amino-acid physicochemical properties, with
\emph{zero} free parameters trained on patient data: every prediction is a
geometric operation reproducible from the algorithm alone, independent of any
training cohort.
I do not present this as the method xMDT-HPB should adopt --- clinical text,
imaging, and laboratory trajectories are different problems with their own
mature approaches.
I present it as evidence that I understand, concretely and from the building
side, the difference between a model that is \emph{explained} and one that is
\emph{explainable}, which is exactly the distinction this project has to get
right for clinicians to trust its output.

\textbf{Hepato-pancreato-biliary, specifically.}
The pancreatic component of HPB is where my existing work already touches the
project's clinical domain.
Syndrome's neoantigen and immune-escape analysis was developed with
KRAS-mutant pancreatic ductal adenocarcinoma as a worked case --- the immune
evasion that makes PDAC so difficult is one of the disease states the framework
was built to characterise.
I am not coming to HPB cancer cold; I have read and modelled in this space,
and I would bring that grounding while learning the radiological and
clinical-narrative modalities that are new to me.

\textbf{Multimodal and longitudinal data.}
The project integrates clinical text, multimodal imaging, and laboratory
trajectories.
Multimodal fusion and longitudinal inference are where my independent work
already lives.
I have built systems that combine physically distinct signal channels into a
single inference, and a framework for monitoring \emph{laboratory trajectories}
as bounded, uncertainty-honest regions rather than point estimates --- inferring
a full state from partial, noisy, longitudinal observation with a provable
convergence guarantee.
That last problem --- reconstructing a trustworthy trajectory from incomplete
clinical measurement --- is structurally the laboratory-data problem xMDT-HPB
faces, and it is one I have worked on directly.

\textbf{Open-source, reproducible, benchmarked --- how I already work.}
The project commits to releasing everything as open-source software and
reusable artifacts, with reproducible experiments and benchmarking pipelines.
This is not an adjustment for me; it is how I build.
My work is public on GitHub, my validation suites are deterministic and
seeded so that results reproduce exactly from a stated seed, and every claim I
make is checkable against an explicit reference rather than asserted.
I would contribute to the project's benchmarking and quality-assurance work
from day one, because reproducible, independently-verifiable results are the
standard I already hold myself to.

\textbf{Why the structured path, and why now.}
I was previously a doctoral candidate in computational lipidomics at TUM, where
I worked on machine learning for large clinical omics cohorts and distributed
inference across federated sites, before funding ended and I left the programme.
Since then I have worked independently and built a great deal --- but the thing
I most want now is precisely what this position provides and independent work
does not: structured supervision, a real interdisciplinary team of clinicians
and engineers, and clinical data and standards (MII, FHIR) I can only learn
inside a group that works with them.
I am genuinely prepared to familiarise myself with medical standards, data
protection, and good scientific practice, and to do the careful, documented,
team-aligned work the project describes.
I would rather ground my work in a real clinical setting under supervision than
continue to extend it alone, and this project is exactly that setting.

\textbf{Background.}
I hold an MSc in Computational Biology (Universität Tübingen) and an MSc in
Pharmaceutical Biotechnology (HAW Hamburg), with a BSc in Biochemistry and Cell
Biology (Jacobs University Bremen).
I conducted doctoral research in Computational Lipidomics at TUM (2019--2021).
I work daily in Python and have hands-on experience with machine-learning
frameworks and the full software-engineering discipline the role asks for:
version control, documentation, testing, and reproducible pipelines.
I have lived and worked in Germany since 2011 and hold permanent residence, so
there is no work-authorisation or sponsorship question.
I completed both of my master's degrees --- Pharmaceutical Biotechnology and
Computational Biology --- entirely in German, and taught master's-level material
in German at TUM.
Two full German-language degrees are, I would suggest, more direct evidence of
the academic and clinical German this role requires than any certificate:
I have already done sustained scientific work in German at exactly this level.
English is native.

I would welcome the chance to discuss how this background fits xMDT-HPB, and I
can start as soon as the project requires.

\vspace{1.5em}
Yours sincerely,

\vspace{2.5em}
\textbf{Kundai Farai Sachikonye}

\end{document}
